% year: תשפ"ה
% type: בוחן
% semester: א
% source: exams/בחנים/תשפ״ה, בוחן, 9141.pdf
% part: א
% number: 3
% points: 20
% categories: functions

\begin{question}
חשבו את הפונקציה ההופכית לפונקציה $f(x)=\ln(1+\sqrt{x})$ ורשמו את תחום ההגדרה של הפונקציה ההפוכה.
\end{question}

\begin{hint}
תחום ההגדרה של הפונקציה ההפוכה הוא התמונה של $f$. מצאו את התמונה, ואז חלצו את $x$ מהמשוואה $y=\ln(1+\sqrt x)$.
\end{hint}

\begin{solution}
\textbf{תחום ההגדרה והתמונה של $f$.} $\sqrt x$ מוגדר עבור $x\ge0$, ואז $1+\sqrt x\ge1>0$ ולכן $\ln(1+\sqrt x)$ מוגדר. לכן $D_f=[0,\infty)$.

$f$ עולה ממש: $\sqrt x$ עולה ממש על $[0,\infty)$ ו-$\ln$ עולה ממש, והרכבה של פונקציות עולות ממש היא עולה ממש. בפרט $f$ חד-חד-ערכית ולכן הפיכה על התמונה שלה.

התמונה: $f(0)=\ln 1=0$ ולכן $f(x)\ge0$ לכל $x\ge0$. לכל $y\ge0$ נמצא בהמשך $x\ge0$ עם $f(x)=y$, ולכן $\operatorname{Im}f=[0,\infty)$.

\textbf{חישוב ההופכית.} יהי $y\ge0$. נפתור $y=\ln(1+\sqrt x)$:
\[
e^y=1+\sqrt x\iff \sqrt x=e^y-1 .
\]
מכיוון ש-$y\ge0$ מתקיים $e^y-1\ge0$, ולכן יש פתרון יחיד $x=(e^y-1)^2\ge0$. (עבור $y<0$ היה מתקבל $\sqrt x<0$ — אין פתרון, כצפוי מכך ש-$y$ אינו בתמונה.)

לכן
\[
\boxed{f^{-1}(x)=\left(e^x-1\right)^2,\qquad D_{f^{-1}}=[0,\infty).}
\]
\textbf{בדיקה:} עבור $x\ge0$: $f(f^{-1}(x))=\ln\left(1+\sqrt{(e^x-1)^2}\right)=\ln\left(1+|e^x-1|\right)=\ln(e^x)=x$, כי $e^x-1\ge0$. שימו לב שהגבלת התחום ל-$[0,\infty)$ הכרחית: הנוסחה $(e^x-1)^2$ מוגדרת לכל $x$, אבל עבור $x<0$ היא אינה הפכית של $f$.
\end{solution}
